\begin{abstract}
The Herculaneum papyri, carbonized by the eruption of Mount Vesuvius in 79~AD, represent the only surviving library from classical antiquity. Recent advances in X-ray computed tomography (CT) and deep learning have enabled virtual unwrapping and ink detection in these fragile scrolls, culminating in the Vesuvius Challenge (2023-2025), which produced the first machine-readable text from intact scrolls. However, the winning solutions exist solely as competition code-lacking theoretical grounding, systematic documentation, and reproducibility as scientific baselines. In particular, the role of depth (Z-axis) layer selection in volumetric ink detection remains poorly understood, with practitioners relying on empirically chosen configurations without principled justification.

This work presents a systematic investigation of Z-layer sampling strategies for ink detection in Herculaneum papyri segments. We reverse-engineer and formalize the Grand Prize winning pipeline, then conduct five controlled experiments: (A)~a layer-count ablation varying the number of input depth slices, (B)~a depth-window position sweep across the papyrus cross-section, (C)~a leave-one-out sensitivity analysis identifying critical individual layers, (D)~a cross-segment generalization study across 13 scroll regions, and (E)~a systematic failure analysis of the worst-performing segments. We evaluate two architecturally distinct models-a TimeSformer treating Z-layers as temporal frames with divided space-time attention, and a 2D U-Net treating layers as input channels-to distinguish architecture-specific from universal depth-sampling effects.

Our results demonstrate that five depth layers centered at the papyrus surface (layers~30-34, spanning ${\sim}40$~\textmu m) achieve competitive ink detection-using $5.2\times$ fewer layers than the competition winner's 26-with diminishing returns beyond this count. Layers~32 and~31, located at and near the papyrus surface center, are the most informative on average across tested segments. A 2D U-Net with EfficientNet-B4 backbone, trained for 80 epochs on pseudo-labels, achieves $F_{0.5} = 0.474$ on a held-out test segment, surpassing the Grand Prize winning TimeSformer ($F_{0.5} = 0.379$) by $25\%$. Training beyond 80 epochs causes overfitting, reducing performance to $F_{0.5} = 0.341$ at 100 epochs-a finding with important implications for training budget allocation on small pseudo-label corpora. We provide concrete, reproducible guidelines for Z-layer configuration that balance detection quality against computational cost, establishing a documented scientific baseline for future ink detection research.
\end{abstract}

\begin{IEEEkeywords}
Herculaneum papyri, ink detection, computed tomography, depth sampling, Z-layer selection, virtual unwrapping, deep learning, TimeSformer, U-Net, Vesuvius Challenge
\end{IEEEkeywords}
